\documentclass{article}
\usepackage[utf8]{inputenc}
\usepackage{amsfonts}
\usepackage{amsmath}
\usepackage{amssymb}
\usepackage[normalem]{ulem}
\usepackage{amsthm}
\usepackage{graphicx}
\usepackage{verbatim}
\usepackage{enumitem}
\usepackage{hyperref}
\usepackage{graphicx}
\usepackage{mathtools}
\usepackage{xcolor}
\usepackage{mathdots}
\usepackage{arydshln}
\usepackage{textcomp}
%\usepackage{bbold}
\usepackage[margin=1in]{geometry}
\hypersetup{
	colorlinks=true,
	linkcolor=blue,
	filecolor=magenta,      
	urlcolor=cyan,
}
\newcommand\numberthis{\addtocounter{equation}{1}\tag{\theequation}}
\newcommand\blank{\hphantom{==}}

\DeclareFontFamily{U}{MnSymbolC}{}
\DeclareSymbolFont{MnSyC}{U}{MnSymbolC}{m}{n}
\DeclareFontShape{U}{MnSymbolC}{m}{n}{
	<-6>  MnSymbolC5
	<6-7>  MnSymbolC6
	<7-8>  MnSymbolC7
	<8-9>  MnSymbolC8
	<9-10> MnSymbolC9
	<10-12> MnSymbolC10
	<12->   MnSymbolC12}{}
\DeclareMathSymbol{\intprod}{\mathbin}{MnSyC}{'270}

\numberwithin{equation}{section}
%\numberwithin{theorem}{section}

%%% THEOREM STYLES %%%
\theoremstyle{plain}
\newtheorem{theorem}{Theorem}[section]
\newtheorem*{theorem*}{Theorem}
\newtheorem{lemma}{Lemma}[section]
\newtheorem*{lemma*}{Lemma}
\newtheorem{proposition}{Proposition}
\newtheorem{corollary}{Corollary}[section]
\newtheorem{conjecture}{Conjecture}

%%% DEFINITION STYLES %%%
\theoremstyle{definition}
\newtheorem{definition}{Definition}[section]
\newtheorem*{definition*}{Definition}
\newtheorem{example}{Example}
\newtheorem*{example*}{Example}
\newtheorem{remark}{Remark}
\newtheorem{remark*}{Remark}
\newtheorem{fact}{Fact}
\newtheorem{problem}{Problem}
\newtheorem*{problem*}{Problem}
\newtheorem*{solution*}{Solution}
\newtheorem*{claim*}{Claim}
\newtheorem{exercise}{Exercise}

%%% FUNCTIONS %%%
\newcommand{\Aut}{\mathrm{Aut}}
\newcommand{\Sym}{\mathrm{Sym}}
\newcommand{\Syl}{\mathrm{Syl}}
\newcommand{\intersection}{\mathop{\bigcap}}
\newcommand{\union}{\mathop{\bigcup}}
\newcommand{\disjointunion}{\mathop{\bigsqcup}}
\newcommand{\evaluated}{\Big|}
\newcommand{\re}{\mathrm{Re}}
\newcommand{\im}{\mathrm{Im}}
\newcommand{\Log}{\mathrm{Log}}
\newcommand{\Res}{\mathrm{Res}}
\newcommand{\Sk}{\mathrm{Sk}}
\newcommand{\disc}{\mathrm{disc}}
\newcommand{\then}{\Longrightarrow}
\newcommand{\Int}{\mathrm{int}}
\newcommand{\bd}{\mathrm{bd}}
\newcommand{\diam}{\mathrm{diam}}
\newcommand{\lcm}{\mathrm{lcm}}
\newcommand{\Tr}{\mathrm{Tr}}
\newcommand{\Gal}{\mathrm{Gal}}
\newcommand{\normal}{\trianglelefteq}
\newcommand{\del}{\partial}
\newcommand{\len}{\mathrm{len}}
\newcommand{\res}{\mathrm{res}}
\newcommand{\esssup}{\mathrm{ess} \, \mathrm{sup}}
\newcommand{\wlim}{\displaystyle\mathop{\mathrm{w}\textrm{-}\mathrm{lim}}}
\newcommand{\vlim}{\displaystyle \mathop{\mathrm{v}\textrm{-}\mathrm{lim}}}
\newcommand{\condE}[2]{\mathbb E\left[ #1 \,\middle|\, #2\right]}
\newcommand{\condP}[2]{\mathbb P\left[ #1 \, \middle| \, #2\right]}
\newcommand{\supp}{\mathrm{supp}}
\newcommand{\Poi}{\mathrm{Poi}}

%%% FONTS %%%
\newcommand{\boldZ}{\boldsymbol{\mathrm{Z}}}
\newcommand{\boldC}{\boldsymbol{\mathrm{C}}}
\newcommand{\boldN}{\boldsymbol{\mathrm{N}}}
\newcommand{\N}{\mathbb N}
\newcommand{\Z}{\mathbb Z}
\newcommand{\Q}{\mathbb Q}
\newcommand{\R}{\mathbb R}
\newcommand{\C}{\mathbb C}
\newcommand{\U}{\mathbb U}
\newcommand{\D}{\mathbb D}
\newcommand{\E}{\mathbb E}
\newcommand{\Sbb}{\mathbb S}
\newcommand{\Kbb}{\mathbb K}
\newcommand{\F}{\mathbb F}
\newcommand{\Pbb}{\mathbb P}
\newcommand{\Ts}{\mathcal T}
\newcommand{\Cs}{\mathcal C}
\newcommand{\As}{\mathcal A}
\newcommand{\Ss}{\mathcal S}
\newcommand{\Ms}{\mathcal M}
\newcommand{\Bs}{\mathcal B}
\newcommand{\Ps}{\mathcal P}
\newcommand{\Ds}{\mathcal D}
\newcommand{\Ls}{\mathcal L}
\newcommand{\Ns}{\mathcal N}
\newcommand{\Rs}{\mathcal R}
\newcommand{\Fs}{\mathcal F}
\newcommand{\Os}{\mathcal O}
\newcommand{\Ks}{\mathcal K}
\newcommand{\Es}{\mathcal E}
\newcommand{\Gs}{\mathcal G}
\newcommand{\Us}{\mathcal U}
\newcommand{\Zs}{\mathcal Z}
\newcommand{\fp}{\mathfrak p}
\newcommand{\fm}{\mathfrak m}
\newcommand{\diff}{\,d}
\newcommand{\id}{\mathrm{id}}
\newcommand{\Span}{\mathrm{span}}
\newcommand{\SL}{\mathrm{SL}}
\newcommand*{\charone}{\text{\usefont{U}{bbold}{m}{n}1}}
\newcommand{\dlim}{\displaystyle\lim}
\newcommand{\dsum}{\displaystyle\sum}
\newcommand{\dint}{\displaystyle\int}
\newcommand{\Var}{\mathbf{Var}}
\newcommand{\Cov}{\mathbf{Cov}}
\newcommand{\Lip}{\mathrm{Lip}}
\newcommand{\sgn}{\mathrm{sgn}}

\newcommand{\done}{\flushright$\diamondsuit$}

\let\oldemptyset\emptyset
\let\emptyset\varnothing
\let\oldphi\phi
\let\phi\varphi
\let\oldepsilon\epsilon
\let\epsilon\varepsilon
\let\oldIm\Im
\let\oldRe\Re
\renewcommand{\Im}{\mathrm{Im}\,}
\renewcommand{\Re}{\mathrm{Re}\,}
\let\oldtilde\tilde
\let\tilde\widetilde
\let\oldhat\hat
\let\hat\widehat

\begin{document}

\begin{center}
\Large{Tutorial Exercises Solutions - Week 5} \\
\vspace{3mm}
\large{4001CMD - Mathematical Skills for Computing Professionals} \\
\vspace{5mm}
\end{center}


The following exercises should be attempted before your next tutorial session this week. Problems with an asterisk (\textbf{*}) are particularly challenging! 

\section*{Addition and Sieve Principles}

\begin{problem}
	Write out a proof (using induction on $n$) of the fact that if $A_1, A_2, \ldots, A_n$ are disjoint finite sets, then \[
	\#\left(A_1 \cup A_2 \cup \cdots \cup A_n\right) = (\#A_1) + (\#A_2) + \cdots + (\#A_n).
	\]
\end{problem}

\begin{solution*}
	\begin{proof}
		For $n=1$, this is just saying $\#A_1 = \#A_1$. For $n=2$, this is the addition principle: $\#(A_1 \cup A_2) = \#A_1 + \#A_2$ if $A_1$ and $A_2$ are disjoint. 
		
		Suppose that for $n=k$, we have: 
		\[
		\#(A_1 \cup \cdots \cup A_k) = \#A_1 + \cdots + \#A_k.
		\]
		Let $A_{k+1}$ be another set, disjoint from the $A_1, \ldots, A_k$. By the $n=2$ case, 
		\begin{align*}
		\#\left(A_1 \cup \cdots \cup A_k \cup A_{k+1}\right) &= \#\left(A_1 \cup \cdots \cup A_k\right)+ \#A_{k+1} \\
		&= \#A_1 + \cdots + \#A_k + \#A_{k+1}. 
		\end{align*}
		So if the addition principle holds for $n=k$, it holds for $n=k+1$. By induction, the claim holds for all $n \geq 2$. 
	\end{proof}
\end{solution*}

\begin{problem}
	In a class of 67 mathematics students, 47 can read French, 35 can read German and 23 can read both languages.
	\begin{enumerate}[label=(\alph*)]
		\item How many can read neither language? 
		\item If furthermore 20 can read Russian, of whom 12 can also read French, 11 read German also and 5 read all three languages, how many cannot read any of the three languages? 
	\end{enumerate}
\end{problem}

\begin{solution*}
	(a) Let $F$ be those students who can read French and $G$ those students who can read German. Also let $X$ be the set of all students in the class. We are told: \[
	\#X = 67, \quad \#F = 47, \#G = 35, \quad \#(F \cap G) = 23
	\]
	By the sieve principle, the number of students who cannot read French or German is:
	\[
	\#(X \setminus(F \cup G)) = \#X - \#F - \#G + \#(F \cap G) = 67 - 47 - 35 + 23 = \boxed{8}.
	\]
	(b) If $R$ is the set of students who can read Russian, we are additionally told: 
	\[
	\#R = 20, \quad \#(F \cap R) = 12, \quad \#(G \cap R) = 11, \quad \#(F \cap G \cap R) = 5.
	\]
	So the number of students who cannot read German, French, or Russian is:
	\begin{align*}
	\#(X \setminus(F \cup G \cup R)) &= \#X - \#F - \#G - \#R + \#(F \cap G) + \#(F \cap R) + \#(G \cap R) - \#(F \cap G \cap R) \\
	&= 67 - 47 - 35 - 20 + 23 + 12 + 11 - 5 = \boxed{6}
	\end{align*}
\end{solution*}

\begin{problem}[\textbf{*}]
	Find the number of ways of arranging the letters A, E, M, O, U, Y in a sequence in such a way that neither of the words ME nor YOU occur. (\emph{Hint:} Apply the sieve principle to the set of all rearrangements of the letters A, E, M, O, U, Y.)
\end{problem}

\begin{solution*}
	Let's say $P$ is the set of all permutations of the letters A, E, M, O, U, and Y. Then $\#P = 6!$. Let $Q \subseteq P$ be the set of permutations for which the word ME appears, and let $R \subseteq P$ be the set of permutations for which the word YOU appears. We want to know the size of $P \setminus(Q \cup R)$. By the sieve principle, $$\#(Q \cup R) = \#Q + \#R - \#(Q \cap R).$$ So we need to find $\#Q$, $\#R$, and $\#(Q \cap R)$. 
	
	For $\#Q$, rhere are 5 different positions in a 6-letter permutation that the word ME could appear: MEXXXX, XMEXXX, XXMEXX, XXXMEX, and XXXXME. In each of these permutations, the remaining four letters can themselves be rearranged in 4! ways. So, $\#Q = 5 \times 4! = 5!$. 
	
	For $\#R$, there are 4 different positions in a 6-letter permutation that YOU can appear in: YOUXXX, XYOUXX, XXYOUX, and XXXYOU. The remaining three letters can be rearranged in 3! ways. So $\#R = 4 \times 3! = 4!$. 
	
	For $\#(Q \cap R)$, we can list out the permutations that contain both ME and YOU: 
	\begin{enumerate}
		\item MEYOUA
		\item MEAYOU
		\item AMEYOU
		\item YOUMEA
		\item YOUAME
		\item AYOUME
	\end{enumerate}
	To be more precise, we can also note that a permutation of A, E, M, O, U, Y that must contain the words ME and YOU is effectively a permutation of the words A, ME, and YOU. There are 3! = 6 such permutations. So $\#(Q \cap R) = 3!$. All of this lets us conclude: 
	\[
	\#(Q \cup R) = \#Q + \#R - \#(Q \cap R) = 5! + 4! - 3! = 6((5 \times 4) + 4 - 1) = 138
	\]
	so there are $6! - 138 = \boxed{582}$ ordered sequences of A, E, M, O, U, Y that do not contain the words ME and YOU. 
\end{solution*}

\section*{Ordered and Unordered Selection}

\begin{problem}
	How many national flags can be constructed from three equal vertical strips, using the colours red, white, blue, and green? (It is assumed that colours can be repeated, and that one vertical edge of the flag is distinguished as the ``flagpole side'', so that flipping the flag around results in a different flag.)
\end{problem}

\begin{solution*}
	The problem is basically asking how many ordered selections of three colors can be made out of four, allowing repetition. There are $\boxed{4^3 = 64 \textrm{ possible flags.}}$
\end{solution*}

\begin{problem}
	Write down all of the subsets of the set $X = \{a,b,c\}$. Use the correspondence between $\Ps(X)$ and $\{0,1\}^3$ discussed in lecture to check that your list is complete. 
\end{problem}

\begin{solution*}
	The correspondence is that the $k$th letter in a word is 0 if the $k$th symbol in $\{a,b,c\}$ is not in the set, and 1 otherwise. So the sets and the corresponding words are: 
	\[
	\boxed{\begin{aligned}
		\emptyset \quad &\longleftrightarrow \quad 000 \\
		\{c\} \quad &\longleftrightarrow \quad 001 \\
		\{b\} \quad &\longleftrightarrow \quad 010 \\
		\{b,c\} \quad &\longleftrightarrow \quad 011 \\
		\{a\} \quad &\longleftrightarrow \quad 100 \\
		\{a,c\} \quad &\longleftrightarrow \quad 101 \\
		\{a,b\} \quad &\longleftrightarrow \quad 110 \\
		\{a,b,c\} \quad &\longleftrightarrow \quad 111
		\end{aligned}}
	\]
	Since this systematically counts up all words of length 3 in the alphabet $\{0,1\}$, this list is complete. 
\end{solution*}

\begin{problem}
	In how many ways can we select a batting order of 11 from a pool of 14 cricketers? 
\end{problem}

\begin{solution*}
	$\displaystyle \frac{14!}{(14-11)!} = \boxed{\frac{14!}{3!} = 14,529,715,200 \textrm{ different batting orders.}}$
\end{solution*}

\begin{problem}
	Evaluate $\binom{16}4$ and $\binom{17}5$. 
\end{problem}

\begin{solution*}
	$\boxed{\displaystyle \binom{16}4 = \frac{16!}{4!12!} = 1820,}$ and $\boxed{ \binom{17}5 = \frac{17!}{5!12!} = 6188.}$
\end{solution*}

%\begin{problem}
%	Explain why the number of words of length $n$ in the alphabet $\{0,1\}$ which contain exactly $r$ zeros is $\binom n r$. 
%\end{problem}

\begin{problem}
	Suppose $X = \{a,b,c\}$. How many selections of 2 elements is are possible from $X$, when those selections are
	\begin{enumerate}[label=(\alph*)]
		\item ordered and repeating? 
		\item ordered and non-repeating?
		\item unordered and repeating? 
		\item unordered and non-repeating? 
	\end{enumerate}
\end{problem}

\begin{solution*}
	\begin{enumerate}[label=(\alph*)]
		\item Ordered and repeating: $3^2 = \boxed 9$. 
		\item Ordered and non-repeating: $\frac{3!}{(3-2)!} = \boxed 6$.
		\item Unordered and repeating: $\binom{3+2-1}{2} = \binom 4 2 = \boxed{6}$. 
		\item Unordered and non-repeating: $\binom{3}{2} = \boxed 3$. 
	\end{enumerate}
\end{solution*}

\begin{problem}
	Keys are made by cutting incisions of various depths in a number of positions on a blank key. If there are eight possible depths, how many positions are required to make one million different keys? (\emph{Hint:} To facilitate the calculation, use the fact that $2^{10}$ is slightly greater than 1000.) 
\end{problem}

\begin{solution*}
	Say there are $p$ positions. If there are 8 possible depths, then this is a question about ordered selection with repetition, and so there are $8^p = 2^{3p}$ possible keys that can be made with $p$ positions. Noting that $2^{10}$ is slightly larger than 1000, we can estimate that we need to make around $2^{20}$ different keys. So we need $3p \geq 20$; we can let $p = 7$. So we need $\boxed{7 \textrm{ positions for the keys.}}$. (Note that if we try doing 6 positions, then we'd have $2^{18} = 262144$ keys, which isn't enough.)
\end{solution*}

\begin{problem}
	How many four-letter words can be made from an alphabet of 10 symbols if there are no restrictions on spelling except that no letter can be used more than once? 
\end{problem}

\begin{solution*}
	This is ordered selection without repetition. We choose 4 items out of 10; there are $\frac{10!}{6!} = 10 \times 9 \times 8 \times 7 = \boxed{5040}$ different such words. 
\end{solution*}

\begin{problem}
	Show that when three indistinguishable dice are thrown, there are 56 possible outcomes. What is the number of outcomes when $n$ indistinguishable dice are thrown?
\end{problem}

\begin{solution*}
	To look at the outcome of throwing three indistinguishable dice, we are looking at a situation of unordered selection with replacement, choosing 3 items out of 6. There are $\binom{6+3-1}{3} = \binom{8}{3} = 56$ possible outcomes. If there are $n$ indistinguishable dice, then there are $\binom{6+n-1}{n}$ possible outcomes. 
\end{solution*}

\begin{problem}[\textbf{*}]
	Show that the number of $n$-tuples $(x_1, \ldots, x_n)$ of non-negative integers which satisfy the equation 
	\[
	x_1 + x_2 + \cdots + x_n = r
	\]
	is $\binom{n+r-1}{r}$. (\emph{Hint:} Suppose that an unordered selection, with repetition, of $r$ objects from a set of $n$ objects contains $x_i$ copies of the $i$th object ($1 \leq i \leq n$).)
\end{problem}

\begin{solution*}
	Let $X = \{(x_1, \ldots, x_n) : x_i \in \N_0, x_1 + \cdots + x_n = r\}$ be the set of $n$-tuples of integers whose sum is $r$. Let $A = \{a_1, \ldots, a_n\}$ be a set of $n$ objects, and let $A^*$ denote the set of $r$-fold unordered selections from $A$ with replacement. Define a function $s : X \to A^*$ in the following way: $s(x_1, \ldots, x_n)$ selects $x_1$ copies of $a_1$, then $x_2$ copies of $a_2$, etc., so that in total, $x_1 + x_2 + \cdots + x_n$ items from $A$ are selected; since $x_1 + \cdots + x_n = r$, the resulting selection is an $r$-fold unordered selection with replacement, so $s$ has codomain $A^*$. 
	
	We claim that $s : X \to A^*$ is a bijection. If $x = (x_1, \ldots, x_n) \in X$ and $y = (y_1, \ldots, y_n) \in X$, and if $s(x) = s(y)$, then for $1 \leq i \leq n$, the selection takes $x_i$ copies of $a_i$ from $A$, and also $y_i$ copies of $a_i$ from $a$; if they're the same selection, this means $x_i = y_i$ for all $i$. So $s$ is injective. And, suppose $\mathbf{a} = (a_{i_1}, \ldots, a_{i_r}) \in A^*$ is an $r$-fold unordered selection with replacement, with $1 \leq i_1 \leq \cdots \leq i_r \leq n$ (e.g., for $r=5$ and $n=6$, we might have $\mathbf a = (a_1, a_1, a_3, a_6, a_6)$). For $1 \leq i \leq n$, let $x_i$ be the number of choices of $a_i$ made in $\mathbf a$ (in the above example, we have $x_1 = 2$, $x_2 = 0$, $x_3 = 1$, $x_4 = 0$, $x_5 = 0$, and $x_6 = 2$). Then $x_1 + \cdots + x_n = r$, since $r$ items are selected in total. By definition of the function $s : X \to A^*$, we have $s(x_1, \ldots, x_n) = \mathbf a$ (in the above example, $s(2,0,1,0,0,2) = (a_1, a_1, a_3, a_6, a_6)$). So $s$ is surjective. 
\end{solution*}

\section*{Binomial Theorem}

\begin{problem}
	Write out the formulas for $(1+x)^8$ and $(1-x)^8$. 
\end{problem}

\begin{solution*}
	We can write out the first eight rows of Pascal's triangle: 
	\[
	\begin{array}{ccccccccccccccccc}
	& & & & & & & & 1 & & & & & & & & \\
	& & & & & & & 1 & & 1 & & & & & & & \\
	& & & & & & 1 & & 2 & & 1 & & & & & &\\
	& & & & & 1 & & 3 & & 3 & & 1 & & & & &\\
	& & & & 1 & & 4 & & 6 & & 4 & & 1 & & & &\\
	& & & 1 & & 5 & & 10 & & 10 & & 5 & & 1 & & &\\
	& & 1 & & 6 & & 15 & & 20 & & 15 & & 6 & & 1 & &\\
	& 1 & & 7 & & 21 & & 35 & & 35 & & 21 & & 7 & & 1 &\\
	1 & & 8 & & 28 & & 56 & & 70 & & 56 & & 28 & & 8 & & 1
	\end{array}
	\]
	So:
	\[
	\boxed{(1+x)^8 = x^8 + 8x^7 + 28x^6 + 56x^5 + 70x^4 + 56x^3 + 28x^2 + 8x + 1}.
	\]
	For $(1-x)^8$, the terms are $1^{8-k}(-x)^{k} = (-1)^k x^k$, so the signs alternate. So: 
	\[
	\boxed{(1-x)^8 = x^8 - 8x^7 + 28x^6 - 56x^5 + 70x^4 - 56x^3 + 28x^2 - 8x + 1}.
	\]
\end{solution*}

\begin{problem}
	Calculate the coefficients of: 
	\begin{enumerate}[label=(\alph*)]
		\item $x^5$ in $(1+x)^{11}$
		\item $a^2b^8$ in $(a+b)^{10}$
		\item $a^6 b^6$ in $(a^2 + b^3)^5$
		\item $x^3$ in $(3+4x)^6$
	\end{enumerate}
\end{problem}

\begin{solution*}
	(a) The term is $\displaystyle \binom{11}5 1^{11-5} x^5 = \frac{11!}{5!6!} x^5 = 462 x^5$, so the coefficient is $\boxed{462}$. 
	
	(b) The term is $\displaystyle \binom{10}2 a^2 b^8 = \frac{10!}{2! 8!} a^2 b^8 = 45a^2 b^8$, so the coefficient is $\boxed{45}$. 
	
	(c) Referring to Pascal's triangle in the previous problem, the expansion is:
	\begin{align*}
	(a^2 + b^3)^5 &= (a^2)^5 + 5(a^2)^4 b^3 + 10(a^2)^3 (b^3)^2 + 10(a^2)^2 (b^3)^3 + 5(a^2)(b^3)^4 + (b^3)^5 \\
	&= a^{10} + 5a^8 b^3 + 10a^6 b^6 + 10a^4 b^9 + 5a^2 b^{12} + b^{15}
	\end{align*}
	so the $a^6 b^6$ coefficient is $\boxed{10}$. 
	
	(d) The $x^3$ term is $\displaystyle \binom{6}{3} 3^3(4x)^3 = 20 \cdot 27 \cdot 64 x^3 = 34560x^3$. So the coefficient is $\boxed{34560}$. 
\end{solution*}

\end{document} 